% understory.tex
% Compile with: xelatex understory.tex (run twice)
\documentclass[10pt, letterpaper]{article}

\usepackage[margin=0.55in, top=0.45in, bottom=0.4in]{geometry}
\usepackage[dvipsnames, table]{xcolor}
\usepackage{tikz}
\usepackage[americanresistors, americaninductors]{circuitikz}
\usepackage{booktabs}
\usepackage{tabularx}
\usepackage{colortbl}
\usepackage{enumitem}
\usepackage{listings}
\usepackage{tcolorbox}
\usepackage{fancyhdr}
\usepackage{fontspec}
\usepackage{hyperref}
\usepackage{parskip}
\usepackage{pdflscape}
\setlength{\parskip}{2pt plus 1pt minus 1pt}

\setmainfont{Latin Modern Roman}
\setmonofont{Latin Modern Mono}

\definecolor{espblue}{HTML}{2563eb}
\definecolor{fetred}{HTML}{dc2626}
\definecolor{sigamber}{HTML}{b45309}
\definecolor{loadcyan}{HTML}{0e7490}
\definecolor{protpurple}{HTML}{7c3aed}
\definecolor{growgreen}{HTML}{15803d}
\definecolor{darkbg}{HTML}{1e293b}
\definecolor{lightbg}{HTML}{f8fafc}
\definecolor{codebg}{HTML}{f1f5f9}
\definecolor{bordergray}{HTML}{cbd5e1}
\definecolor{textgray}{HTML}{64748b}
\definecolor{darktext}{HTML}{1e293b}
\definecolor{gndcolor}{HTML}{374151}
\definecolor{vplusred}{HTML}{b91c1c}
\definecolor{divider}{HTML}{94a3b8}

\hypersetup{colorlinks=true, linkcolor=espblue, urlcolor=espblue,
  pdfauthor={AJ}, pdftitle={Herb Garden Light Controller}}

\tcbuselibrary{skins, breakable}
\newtcolorbox{calloutbox}[1][]{%
  colback=blue!3, colframe=espblue!40,
  fonttitle=\bfseries\small, left=5pt, right=5pt,
  top=3pt, bottom=3pt, boxrule=0.5pt, arc=3pt, #1}
\newtcolorbox{codebox}[1][]{%
  colback=codebg, colframe=bordergray,
  left=5pt, right=5pt, top=3pt, bottom=3pt,
  boxrule=0.4pt, arc=2pt, #1}

\lstdefinestyle{circuitpy}{
  language=Python, basicstyle=\ttfamily\tiny,
  keywordstyle=\color{espblue}\bfseries, stringstyle=\color{growgreen},
  commentstyle=\color{textgray}\itshape, backgroundcolor=\color{codebg},
  frame=single, rulecolor=\color{bordergray}, framesep=4pt,
  xleftmargin=4pt, xrightmargin=4pt, breaklines=true,
  showstringspaces=false, tabsize=4,
  morekeywords={digitalio, wifi, socketpool, adafruit_ntp, Server, Request, Response}}

\pagestyle{fancy}
\fancyhf{}
\renewcommand{\headrulewidth}{0pt}
\fancyfoot[L]{\scriptsize\color{textgray} Herb Garden Light Controller · Rev~1 · Feb 2026}
\fancyfoot[C]{\scriptsize\color{textgray} \thepage/2}
\fancyfoot[R]{\scriptsize\color{textgray} CircuitPython · ESP32-S3 · Open Source}

\usepackage{titlesec}
\titleformat{\section}{\large\bfseries\color{darktext}}{}{0em}{}[\vspace{-4pt}\color{bordergray}\rule{\linewidth}{0.4pt}]
\titlespacing*{\section}{0pt}{6pt}{2pt}
\titleformat{\subsection}{\normalsize\bfseries\color{darktext}}{}{0em}{}
\titlespacing*{\subsection}{0pt}{5pt}{1pt}

\setlist[itemize]{nosep, leftmargin=1.2em, label=\textcolor{espblue}{\textbullet}}
\setlist[enumerate]{nosep, leftmargin=1.5em}

\begin{document}

% ══════════════════════════════════════════════════════════════
% PAGE 1: LANDSCAPE — Schematic + Notes + BOM
% ══════════════════════════════════════════════════════════════

\begin{landscape}
\thispagestyle{fancy}

{\huge\bfseries\color{darktext} Herb Garden Light Controller}\hfill
{\color{textgray}\small ESP32-S3 + MOSFET · CircuitPython · WiFi Scheduler}

\vspace{1pt}

\section{Schematic --- Low-Side N-Channel MOSFET Switch}

\vspace{-4pt}

\begin{center}
\begin{circuitikz}[
  american,
  scale=0.82, transform shape,
  thick,
]

% ══════════════════════════════════════════════════
% LAYOUT GRID (conceptual, not drawn):
% x: 0-4 digital | 5 boundary | 6-16 analog | 17-20 load
% y: main signal path at y=3
%    components extend y=-3 to y=8
% ══════════════════════════════════════════════════

% ══════════════════════════════════════════════════
% DOMAIN BOUNDARY — vertical dashed line
% This is the conceptual division
% ══════════════════════════════════════════════════

\draw[divider, thick, dash pattern=on 8pt off 4pt]
  (5.0, 8.5) -- (5.0, -3.8);
\node[font=\small\itshape, color=divider, rotate=90, anchor=south]
  at (4.6, 2.5) {software $\rightarrow$ electricity};

% Domain labels at top
\node[font=\normalsize\bfseries, color=espblue!60] at (2.5, 8.2) {DIGITAL};
\node[font=\scriptsize, color=textgray] at (2.5, 7.6) {CircuitPython on ESP32};
\node[font=\scriptsize, color=textgray] at (2.5, 7.1) {\texttt{light.value = True}};

\node[font=\normalsize\bfseries, color=fetred!80] at (11.0, 8.2) {ANALOG};
\node[font=\scriptsize, color=textgray] at (11.0, 7.6) {Real components, real current};

\node[font=\normalsize\bfseries, color=growgreen!70] at (19.0, 8.2) {LOAD};
\node[font=\scriptsize, color=textgray] at (19.0, 7.6) {Existing grow light};

% ══════════════════════════════════════════════════
% U1: ESP32 — small, left side, de-emphasized
% ══════════════════════════════════════════════════

\draw[thin, fill=blue!3, draw=espblue!40, rounded corners=4pt]
  (0.2, 1.5) rectangle (3.8, 5.0);

\node[font=\small\bfseries, color=espblue!50] at (2.0, 4.5) {U1};
\node[font=\small, color=espblue!50] at (2.0, 3.9) {QT Py ESP32-S3};
\node[font=\scriptsize, color=textgray!70] at (2.0, 3.3) {Adafruit 5700};
\node[font=\scriptsize, color=textgray!70] at (2.0, 2.7) {WiFi · NTP · Web UI};

% A0 pin — this is the output that matters
\fill[sigamber] (3.8, 3.0) circle (3pt);
\node[font=\scriptsize\bfseries, color=sigamber, anchor=east] at (3.6, 3.0) {A0};

% GND pin
\fill[gndcolor!50] (3.8, 2.0) circle (2pt);
\node[font=\scriptsize, color=gndcolor!50, anchor=east] at (3.6, 2.0) {GND};

% USB-C — tiny, above
\draw[very thin, fill=gray!3, draw=textgray!30, rounded corners=2pt]
  (0.8, 5.5) rectangle (3.2, 6.2);
\node[font=\scriptsize, color=textgray!50] at (2.0, 5.85) {USB-C 5\,V};
\draw[textgray!30, very thin, dashed] (2.0, 5.5) -- (2.0, 5.0);

% ══════════════════════════════════════════════════
% SIGNAL WIRE: A0 across boundary to junction
% Thickens as it crosses into analog domain
% ══════════════════════════════════════════════════

\draw[sigamber, thick] (3.8, 3.0) -- (5.0, 3.0);
\draw[sigamber, line width=2.5pt] (5.0, 3.0) -- (8.0, 3.0);

% Label on wire
\node[font=\normalsize, color=sigamber, anchor=south]
  at (6.5, 3.2) {3.3\,V when ON};

% ══════════════════════════════════════════════════
% JUNCTION — clearly visible
% ══════════════════════════════════════════════════

\fill[sigamber] (8.0, 3.0) circle (5pt);

% ══════════════════════════════════════════════════
% R1: 10kΩ PULL-DOWN
% Below junction, generous vertical space
% Label to the LEFT, well clear of Q1
% ══════════════════════════════════════════════════

\draw[line width=2pt, color=darktext]
  (8.0, 3.0) -- (8.0, 2.2)
  to[R, color=darktext, l_={\large\bfseries 10\,k$\Omega$}] (8.0, -1.0)
  -- (8.0, -2.0) node[ground, color=gndcolor, scale=1.2]{};

% R1 ref — LEFT side, clear space
\node[font=\Large\bfseries, color=darktext, anchor=east] at (7.1, 0.6) {R1};
\node[font=\small, color=textgray, anchor=east] at (7.1, 0.0) {pull-down};

% ══════════════════════════════════════════════════
% GATE WIRE: junction to Q1
% ══════════════════════════════════════════════════

\draw[sigamber, line width=2.5pt] (8.0, 3.0) -- (11.5, 3.0);

% ══════════════════════════════════════════════════
% Q1: MOSFET — THE HERO
% Centered at x=12, generous clearance all around
% Labels BELOW and to the RIGHT only
% ══════════════════════════════════════════════════

\draw (12.0, 3.0) node[nigfete, anchor=gate,
  line width=2pt, color=fetred, scale=1.5](Q1){};

% Pin labels — on the pins, spaced away
\node[font=\Large\bfseries, color=sigamber, anchor=east] at (11.6, 3.0) {G};
\node[font=\Large\bfseries, color=loadcyan] at (12.8, 6.8) {D};
\node[font=\Large\bfseries, color=gndcolor] at (12.8, -0.8) {S};

% Q1 ref designator — RIGHT side, well below drain
\node[font=\Huge\bfseries, color=fetred, anchor=west] at (14.0, 4.5) {Q1};

% Part number — BELOW ref, still right side
\node[font=\large\bfseries, color=fetred, anchor=west] at (14.0, 3.6) {STP55NF06L};
\node[font=\normalsize, color=textgray, anchor=west] at (14.0, 2.8) {N-Ch MOSFET};
\node[font=\normalsize, color=textgray, anchor=west] at (14.0, 2.1) {60\,V · 55\,A · TO-220};

% Source to GND
\draw[line width=2pt] (Q1.source) -- ++(0, -1.2)
  node[ground, color=gndcolor, scale=1.2]{};

% ══════════════════════════════════════════════════
% D1: FLYBACK DIODE
% To the RIGHT of Q1, own vertical space
% No overlap with Q1 labels
% ══════════════════════════════════════════════════

% Route: Source right, up through diode, left to drain
\draw[protpurple, line width=1.8pt]
  (Q1.source) -- ++(3.5, 0) coordinate(D1bot);

\draw[protpurple, line width=1.8pt]
  (D1bot) to[D, color=protpurple, l={\Large\bfseries D1}] ++(0, 5.0)
  coordinate(D1top);

\draw[protpurple, line width=1.8pt]
  (D1top) -- (Q1.drain);

% D1 part number — to the RIGHT of diode, clear space
\node[font=\large\bfseries, color=protpurple, anchor=west] at (17.0, 3.0) {1N4007};
\node[font=\small, color=textgray, anchor=west] at (17.0, 2.3) {flyback};
\node[font=\small, color=textgray, anchor=west] at (17.0, 1.7) {protection};

% ══════════════════════════════════════════════════
% DRAIN WIRE → LOAD
% ══════════════════════════════════════════════════

\draw[loadcyan, line width=2.5pt]
  (Q1.drain) -- ++(0, 0.8) -- ++(5.5, 0) coordinate(loadwire);

% ══════════════════════════════════════════════════
% LOAD SIDE: Grow Light + PSU — subdued, far right
% ══════════════════════════════════════════════════

% Grow Light — small box, upper right
\draw[thin, fill=green!3, draw=growgreen!40, rounded corners=4pt]
  (17.8, 5.8) rectangle (21.0, 7.5);
\node[font=\normalsize\bfseries, color=growgreen!60] at (19.4, 7.0) {Grow Light};
\node[font=\scriptsize, color=textgray!70] at (19.4, 6.3) {LED Panel + Driver};

% Wire from drain to grow light GND
\draw[loadcyan, line width=2.5pt] (loadwire) -- ++(0.3,0) -- ++(0, 0.6);
\node[font=\small\bfseries, color=loadcyan, fill=white, inner sep=1pt,
  anchor=east] at (18.3, 6.5) {LED GND};

% V+ rail (right edge)
\draw[vplusred, line width=2.5pt] (20.5, 0.3) -- (20.5, 5.8);
\node[font=\normalsize\bfseries, color=vplusred, rotate=90, anchor=south]
  at (20.9, 3.0) {V+};

% V+ connects to grow light
\draw[vplusred, line width=2pt] (20.5, 5.8) -- (20.5, 5.8);

% LED PSU — small box, lower right
\draw[thin, fill=red!2, draw=vplusred!40, rounded corners=4pt]
  (17.8, -1.2) rectangle (21.0, 0.3);
\node[font=\normalsize\bfseries, color=vplusred!60] at (19.4, -0.1) {12\,V / 24\,V DC};
\node[font=\scriptsize, color=textgray!70] at (19.4, -0.7) {LED Power Supply};

% ══════════════════════════════════════════════════
% COMMON GND BUS
% ══════════════════════════════════════════════════

\draw[gndcolor, thick, dashed]
  (3.8, 2.0) -- (4.8, 2.0) -- (4.8, -3.2) -- (19.0, -3.2) -- (19.0, -1.2);
\node[font=\normalsize, color=gndcolor, fill=white, inner sep=2pt]
  at (11.5, -3.2) {Common GND Bus};

% ══════════════════════════════════════════════════
% Legend — single row, bottom
% ══════════════════════════════════════════════════

\draw[thin, fill=white, draw=bordergray, rounded corners=2pt]
  (0.2, -4.4) rectangle (21.0, -3.65);
\node[font=\small, anchor=west] at (0.5, -4.0) {%
  \textcolor{vplusred}{\Large$\bullet$}\;\textcolor{darktext}{V+ Power}\qquad
  \textcolor{sigamber}{\Large$\bullet$}\;\textcolor{darktext}{GPIO Signal}\qquad
  \textcolor{loadcyan}{\Large$\bullet$}\;\textcolor{darktext}{Load Path}\qquad
  \textcolor{gndcolor}{\Large$\bullet$}\;\textcolor{darktext}{Ground}\qquad
  \textcolor{protpurple}{\Large$\bullet$}\;\textcolor{darktext}{Protection}\qquad
  \textcolor{gndcolor}{- - -}\;\textcolor{darktext}{GND Bus}\qquad
  \textcolor{divider}{- - -}\;\textcolor{darktext}{Domain Boundary}};

% Title block (inline with legend)
\node[font=\small, color=textgray, anchor=east] at (20.8, -4.0)
  {Sheet 1/1 · Rev 1 · Feb 2026};

\end{circuitikz}
\end{center}

\vspace{-6pt}

% ══════════════════════════════════════════════════
% Two columns: How it works + BOM
% ══════════════════════════════════════════════════

\noindent\begin{minipage}[t]{0.48\linewidth}
\subsection{How It Works}

\begin{enumerate}[leftmargin=1.5em, itemsep=1pt, label=\textbf{\arabic*.}]
\item \textbf{U1} (digital) connects to WiFi, syncs clock via NTP.
\item Schedule says ON $\rightarrow$ CircuitPython sets \texttt{A0 = True}.
\item Pin A0 outputs 3.3\,V. \textit{Software becomes electricity.}
\item 3.3\,V on \textbf{Q1} gate (analog) opens the MOSFET channel.
\item Current: Grow Light $\rightarrow$ \textbf{Q1} D$\rightarrow$S $\rightarrow$ GND. Light on.
\item Schedule says OFF $\rightarrow$ A0 LOW $\rightarrow$ \textbf{Q1} closes. Light off.
\end{enumerate}

\vspace{2pt}
{\small\color{textgray}
\textbf{R1:} Keeps gate LOW during ESP32 boot (prevents light flash).\\
\textbf{D1:} Clamps voltage spikes when Q1 switches off.}

\end{minipage}\hfill
\begin{minipage}[t]{0.50\linewidth}

\subsection{Bill of Materials --- DigiKey}

\renewcommand{\arraystretch}{1.1}
{\small
\begin{tabularx}{\linewidth}{c c l >{\ttfamily\scriptsize}l c r}
\toprule
\rowcolor{darkbg}
\textcolor{white}{\bfseries \#} &
\textcolor{white}{\bfseries Ref} &
\textcolor{white}{\bfseries Component} &
\textcolor{white}{\bfseries\rmfamily DigiKey \#} &
\textcolor{white}{\bfseries Qty} &
\textcolor{white}{\bfseries Price} \\
\midrule
\rowcolor{lightbg}
1 & U1 & QT Py ESP32-S3 & 1528-5700-ND & 1 & \$12.50 \\
2 & Q1 & STP55NF06L & 497-6742-5-ND & 2 & \$2.07 \\
\rowcolor{lightbg}
3 & D1 & 1N4007 Diode & 1N4007DICT-ND & 1 & \$0.10 \\
4 & R1 & 10\,k$\Omega$ \textrm{1/4W} & CF14JT10K0CT-ND & 5 & \$0.10 \\
\rowcolor{lightbg}
5 & -- & Jumper wires & \textrm{Global Spec.} & 1 & \$5.60 \\
6 & -- & Breadboard & \textrm{(search)} & 1 & $\sim$\$4 \\
\bottomrule
\end{tabularx}
}

\vspace{3pt}
{\scriptsize\color{textgray}
\textbf{Total:} $\sim$\$28 \enspace$\cdot$\enspace
\textbf{Order:} \texttt{digikey.com/en/mylists}\\
\textbf{Before starting:} Open grow light base, find DC voltage (12\,V or 24\,V).}

\end{minipage}

\end{landscape}

% ══════════════════════════════════════════════════════════════
% PAGE 2: PORTRAIT — Software
% ══════════════════════════════════════════════════════════════

\section{Software --- CircuitPython}

The ESP32 runs \textbf{CircuitPython}---Python that executes directly on the microcontroller.
No compiler. No IDE. Plug in via USB-C, a flash drive appears, drag files onto it.
Edit code, save, it runs.

\begin{calloutbox}[title=Why this is cool]
The ESP32 runs a tiny web server. Open a browser on your phone, visit
\texttt{http://herbgarden.local}, and you get a scheduling UI.
No app. No cloud. All running on a chip the size of a postage stamp under the kitchen cabinet.
\end{calloutbox}

\vspace{-2pt}

\subsection{File Structure on the ESP32}

\begin{codebox}
\ttfamily\small
/code.py\hfill\textrm{\scriptsize\color{textgray} Main program---runs on boot}\\
/lib/\\
\quad adafruit\_httpserver/\hfill\textrm{\scriptsize\color{textgray} Web server library}\\
\quad adafruit\_ntp.mpy\hfill\textrm{\scriptsize\color{textgray} Network time sync}\\
/static/\\
\quad index.html\hfill\textrm{\scriptsize\color{textgray} Schedule UI (what your phone sees)}\\
/schedule.json\hfill\textrm{\scriptsize\color{textgray} Saved on/off times}\\
/settings.toml\hfill\textrm{\scriptsize\color{textgray} WiFi creds, timezone, quiet hours}
\end{codebox}

\vspace{-2pt}

\subsection{The Main Loop (Simplified)}

\begin{lstlisting}[style=circuitpy]
import time, board, digitalio, json, wifi, socketpool
import adafruit_ntp
from adafruit_httpserver import Server, Request, Response

# GPIO pin that controls the MOSFET gate
light = digitalio.DigitalInOut(board.A0)
light.direction = digitalio.Direction.OUTPUT

# Connect to WiFi (creds from settings.toml)
pool = socketpool.SocketPool(wifi.radio)
ntp = adafruit_ntp.NTP(pool, tz_offset=-8)  # Pacific Time

# Load saved schedule
with open("/schedule.json", "r") as f:
    schedule = json.load(f)

server = Server(pool)

@server.route("/")
def index(request: Request):
    with open("/static/index.html", "r") as f:
        return Response(request, f.read(), content_type="text/html")

server.start(str(wifi.radio.ipv4_address))

while True:
    server.poll()                              # handle web requests
    now = ntp.datetime
    light.value = should_be_on(now, schedule)  # True/False -> HIGH/LOW
    time.sleep(1)
\end{lstlisting}

\subsection{Build Phases}

\begin{description}[leftmargin=0em, font=\bfseries\small, nosep, itemsep=2pt]
\item[Phase 1 --- Prove it works {\normalfont\color{textgray}(1 afternoon)}]
  Flash CircuitPython. Connect WiFi. Sync time. Toggle onboard LED on schedule. Web UI on phone. No MOSFET yet.

\item[Phase 2 --- Real load {\normalfont\color{textgray}(1 evening)}]
  Wire MOSFET on breadboard. Connect to grow light's DC ground. Control actual light from web UI.

\item[Phase 3 --- Install {\normalfont\color{textgray}(weekend)}]
  Mount under cabinet. Solder permanent connections. Set schedule. Forget about it.

\item[Phase 4 --- Show off {\normalfont\color{textgray}(optional)}]
  GitHub. Blog post. MQTT for Home Assistant. Soil moisture sensor. Go wild.
\end{description}

\vspace{4pt}

\begin{calloutbox}[title=First step when parts arrive]
Plug QT Py into USB-C. A drive called \texttt{CIRCUITPY} appears.
If it doesn't, visit \texttt{circuitpython.org/downloads} and flash the latest build
for ``Adafruit QT Py ESP32-S3.'' Then create \texttt{settings.toml} with your WiFi credentials.
\end{calloutbox}

\vfill
\begin{center}
{\scriptsize\color{textgray}
Hardware: Adafruit + STMicro \quad$\cdot$\quad
Software: CircuitPython \quad$\cdot$\quad
Basil: indifferent to all of this}
\end{center}

\end{document}
